% Part: incompleteness
% Chapter: incompleteness-provability
% Section: fixed-point-lemma

\documentclass[../../../include/open-logic-section]{subfiles}

\begin{document}

\olfileid[fr]{inc}{inp}{fix}

\olsection{Le lemme du point fixe}

\begin{explain}
Le lemme du point fixe affirme que, pour toute
!!{formula}~$!B(x)$ à une seule variable libre, il existe un !!{sentence}~$!A$
tel que $\Th{T}\Proves !A\liff !B(\gn{!A})$,
pourvu que $\Th{T}$ étende~$\Th{Q}$. Dans le cas de
la phrase du menteur, nous voudrions que $!A$ soit
équivalente, de façon démontrable dans~$\Th{T}$,
à « $\gn{!A}$ est fausse », autrement dit à l'énoncé
selon lequel $\Gn{!A}$ est le numéro de Gödel d'un
!!{sentence} faux. Pour saisir l'idée de la preuve,
comparons-la à la paraphrase informelle de~$!A$
donnée par Quine : « donne une phrase fausse quand on
le fait précéder de sa propre citation » donne une
phrase fausse quand on le fait précéder de sa propre
citation. On peut appeler \emph{diagonalisation}
l'opération qui consiste à placer devant une
expression sa propre citation pour former un
!!{sentence}, et appeler le résultat sa diagonalisée.
Ainsi, la diagonalisée de « donne une phrase fausse
quand on le fait précéder de sa propre citation »
est précisément la phrase de Quine. Remarquons que
celle-ci n'est pas la diagonalisée de « donne une
phrase fausse », mais de « donne une phrase fausse
quand on le fait précéder de sa propre citation ».
La propriété que l'on diagonalise fait donc
elle-même intervenir la diagonalisation.

Dans le langage de l'arithmétique, nous formons
la citation d'une !!{formula} à une variable libre
en calculant son numéro de Gödel, puis en substituant
à la variable libre le numéral de ce nombre. La
diagonalisée de~$!E(x)$ est $!E(\num{n})$, où
$n=\Gn{!E(x)}$. Désormais, abrégeons
$\num{\Gn{!E(x)}}$ en $\gn{!E(x)}$. Si $!B(x)$
exprime « $x$ code un !!{sentence} faux », alors
« donne une phrase fausse quand on le fait précéder
de sa propre citation » revient à « donne un
!!{sentence} faux quand on lui applique le numéro
de Gödel de sa diagonalisée ». Supposons que nous
disposions d'un symbole~$\Obj{diag}$ pour la fonction
$\fn{diag}(n)$ qui calcule le numéro de Gödel de la
diagonalisée de la !!{formula} de numéro~$n$.
Nous pourrions écrire $!E(x)$ sous la forme
$!B(\Obj{diag}(x))$. La phrase du menteur de Quine
serait alors sa diagonalisée, à savoir
$!E(\gn{!E(x)})$, ou encore
$!B(\Obj{diag}(\gn{!B(\Obj{diag}(x))}))$.
La construction fonctionnerait pour toute autre
propriété~$!B(x)$. Pour le théorème d'incomplétude,
nous choisirons $!B(x)$ pour exprimer « $x$ n'est
pas !!{derivable} dans~$\Th{T}$ ». Alors $!E(x)$
exprimera « donne un !!{sentence} non
!!{derivable} dans~$\Th{T}$ lorsqu'on lui applique
le numéro de Gödel de sa diagonalisée ».

Pour formaliser ce raisonnement dans~$\Th{T}$,
il faut représenter la fonction~$\fn{diag}$.
Elle est calculable, et même récursive primitive :
si $n$ code une !!{formula}~$!E(x)$ dont les variables libres
sont parmi~$x$, alors $\fn{diag}(n)$ renvoie
le code de~$!E(\gn{!E(x)})$; sinon, posons $\fn{diag}(n)=0$.
Rappelons que $\gn{!E(x)}$ est le numéral standard
du nombre $\Gn{!E(x)}$, soit $\num{\Gn{!E(x)}}$.
Si $\Obj{diag}$ était un symbole de fonction de
$\Th{T}$ représentant $\fn{diag}$, nous pourrions
prendre pour $!A$ la !!{formula}
$!B(\Obj{diag}(\gn{!B(\Obj{diag}(x))}))$.
Observons que
\begin{align*}
\fn{diag}(\Gn{!B(\Obj{diag}(x))}) & =
\Gn{!B(\Obj{diag}(\gn{!B(\Obj{diag}(x))}))} \\
& = \Gn{!A}.
\end{align*}
Si $\Th{T}$ peut démontrer
\[
\Obj{diag}(\gn{!B(\Obj{diag}(x))})=\gn{!A},
\]
elle peut aussi démontrer
$!B(\Obj{diag}(\gn{!B(\Obj{diag}(x))}))
\liff !B(\gn{!A})$. Or, par définition, le membre
de gauche est~$!A$.

Bien entendu, $\Obj{diag}$ n'est généralement pas
un symbole de fonction du langage de~$\Th{T}$,
et certainement pas de celui de~$\Th{Q}$.
Mais $\fn{diag}$ étant calculable, une
!!{formula}~$!D_{\fn{diag}}(x,y)$ la représente
dans~$\Th{Q}$. Au lieu d'écrire
$!B(\Obj{diag}(x))$, nous pouvons donc écrire
$\lexists[y][(!D_{\fn{diag}}(x,y)\land !B(y))]$.
Le reste de la preuve esquissée fonctionne alors,
et fonctionne déjà dans~$\Th{Q}$.
\end{explain}

\begin{lem}
\ollabel{lem:fixed-point}
Soit $!B(x)$ une !!{formula} dont la seule variable
libre est~$x$. Il existe un !!{sentence}~$!A$
tel que $\Th{Q}\Proves !A\liff !B(\gn{!A})$.
\end{lem}

\begin{proof}
Étant donnée $!B(x)$, posons
$!E(x)=\lexists[y][(!D_{\fn{diag}}(x,y)\land !B(y))]$,\footnote{La variable~$y$ est distincte de~$x$. Les variables liées sont renommées au préalable pour que $!B(y)$ et les substitutions soient sans capture.}
et prenons pour~$!A$ sa diagonalisée, c'est-à-dire
$!E(\gn{!E(x)})$.

Puisque $!D_{\fn{diag}}$ représente $\fn{diag}$ et
que $\fn{diag}(\Gn{!E(x)})=\Gn{!A}$, $\Th{Q}$
peut démontrer
\begin{align}
  & !D_{\fn{diag}}(\gn{!E(x)},\gn{!A})
    \ollabel{repdiag1} \\
  & \lforall[y][(!D_{\fn{diag}}(\gn{!E(x)},y)
    \lif\eq[y][\gn{!A}])]. \ollabel{repdiag2}
\end{align}
Montrons que $\Th{Q}\Proves !A\liff !B(\gn{!A})$.
Nous raisonnerons informellement, en n'utilisant que
la logique et les faits démontrables dans~$\Th{Q}$.

Supposons d'abord~$!A$, soit $!E(\gn{!E(x)})$.
Par définition de~$!E(x)$, cet énoncé est
\[
\lexists[y][(!D_{\fn{diag}}(\gn{!E(x)},y)
\land !B(y))].
\]
Prenons un tel~$y$. La condition d'unicité de la
représentation donne $y=\gn{!A}$, comme l'indique
la seconde formule ci-dessus (\olref{repdiag2}).
De $!B(y)$, nous déduisons donc $!B(\gn{!A})$.

Supposons maintenant $!B(\gn{!A})$. La première
formule ci-dessus (\olref{repdiag1}) donne
\begin{align*}
& !D_{\fn{diag}}(\gn{!E(x)},\gn{!A})
  \land !B(\gn{!A}).
\intertext{Il s'ensuit que}
& \lexists[y][(!D_{\fn{diag}}(\gn{!E(x)},y)
  \land !B(y))].
\end{align*}
Mais c'est $!E(\gn{!E(x)})$, c'est-à-dire~$!A$.
\end{proof}

\begin{digress}
Comparez cette preuve à celle du lemme du point
fixe en théorie de la calculabilité. Ici, nous
définissons un \emph{énoncé} à partir de lui-même,
alors que, là, nous définissions une \emph{fonction}
à partir d'elle-même. Hormis cette différence,
l'idée est la même.
\end{digress}

\begin{prob}
Une !!{formula}~$!A(x)$ est une \emph{définition
de la vérité} si, pour tout !!{sentence}~$!B$,
$\Th{Q}\Proves !B\liff !A(\gn{!B})$.
Montrez, à l'aide du lemme du point fixe, qu'aucune
!!{formula} n'est une définition de la vérité.
\end{prob}

\end{document}
